% Part overview, set before the part's first chapter.
\checkedwith{9.5.3-dev}{a9b2d49c8}

Most computational work in organic chemistry asks one of a handful of questions: what a molecule looks like, which of its shapes it actually adopts, how stable it is relative to an isomer or a product, how fast it reacts and by which path, and what its spectra should look like. This part, together with the reaction and spectroscopy chapters it points to, shows how each of these is done in ECCE and how the answer from one code is handed to the next.

The first question is about bonding rather than energies. When we draw ethene with two sp$^2$ carbons and a $\pi$ bond, we are using a localised picture that the canonical orbitals of a calculation do not show directly. ECCE-QM can transform the occupied orbitals into localised ones and label each bond with its hybrid composition, so that the textbook picture can be checked against the calculation and its limits seen; Chapter~\ref{ch:localised-bonds} does this for ethane, ethene and ethyne.

Conformers come next, because nearly every property of a flexible molecule is an average over the shapes it adopts at room temperature, and a structure optimised from one drawn geometry is only one of them, and not necessarily the lowest. In ECCE the route is CREST, which searches conformational space with the semi-empirical GFN2-xTB method and returns an ensemble with Boltzmann weights (Chapter~\ref{ch:crest}). The same program finds the likely protonation sites and tautomers of a molecule (Chapter~\ref{ch:protonation}), which matters whenever an acid--base equilibrium or a hydrogen shift decides what is actually present in solution.

These searches are affordable only because xTB is cheap, and Chapter~\ref{ch:xtb} shows what xTB is good for on its own: geometries, a first look at reaction energies, and screening a set of molecules before anything expensive is run. The pattern that runs through the rest of the book starts here. A cheap method ranks many candidates, a density functional calculation in ORCA, Gaussian~16 or ECCE-QM refines the few that matter, and the hand-off between them is made in ECCE, usually by hand with \menu{New from This}, or with a workflow when the chain is long enough to be worth automating (Chapters~\ref{ch:coupling} and \ref{ch:workflows}). Note that the energies of the two levels are never compared with each other; the cheap method only decides which structures go on.

Reactions are treated in the Physical chemistry part, since the methods are not specific to organic chemistry. Reaction and activation energies, enthalpies and Gibbs energies are covered in Chapter~\ref{ch:thermochemistry}; finding a transition state, confirming it by its single imaginary frequency and following the reaction path down to reactants and products with an IRC in Chapter~\ref{ch:reaction-rates}; and rotational barriers and other relaxed scans in Chapter~\ref{ch:scans}. The organic examples that ship with ECCE, substitution against elimination for fluoride and chloroethane, and keto--enol tautomerism with and without a water molecule as a proton relay, were made for these chapters.

Spectra are where the calculations meet experiment. IR and Raman spectra, with the symmetry and composition of each mode, are in Chapter~\ref{ch:vibrations}, the comparison of a calculated spectrum with a measured one in Chapter~\ref{ch:ir-experiment}, UV--vis spectra in Chapter~\ref{ch:excited-states}, and NMR chemical shifts in Chapter~\ref{ch:nmr}. Finally, when the question is about a series of molecules rather than one, the descriptors of Chapter~\ref{ch:qsar} turn a folder of calculations into a table that can be fitted against a measured activity, and Chapter~\ref{ch:docking} shows how a molecule optimised here is docked into a protein.
